\documentclass[a4paper,12pt]{article}
\usepackage[T2A]{fontenc}
\usepackage[utf8]{inputenc}
\usepackage[russian]{babel}
\usepackage{amsmath,amssymb}
\usepackage{geometry}
\geometry{left=2cm,right=2cm,top=2cm,bottom=2cm}
\title{Подробный конспект ответов на экзаменационные вопросы\\
(по учебнику Таненбаума \& Вудхалла, 3-е издание)}
\date{}
\begin{document}
\maketitle

\section*{1. Структура и функции операционных систем. Аппаратная поддержка}

\subsection{Функции ОС}
ОС выполняет две взаимодополняющие функции (глава 1):
\begin{enumerate}
    \item \textbf{Расширенная машина (виртуальная машина):} предоставляет прикладным программам простой, удобный интерфейс, скрывая низкоуровневые детали аппаратуры. Например, вместо программирования регистров контроллера диска (установка номера цилиндра, головки, сектора) ОС даёт абстракцию «файл», который можно открыть, прочитать, записать и закрыть.
    \item \textbf{Менеджер ресурсов:} управляет распределением ресурсов (процессорного времени, памяти, устройств ввода-вывода) между конкурирующими процессами. Ресурсы могут мультиплексироваться двумя способами: 
    \begin{itemize}
        \item \textbf{Во времени:} процессы по очереди используют один ресурс (например, CPU).
        \item \textbf{В пространстве:} каждый процесс получает часть ресурса (например, раздел памяти).
    \end{itemize}
\end{enumerate}

\subsection{Система прерываний}
\textbf{Прерывание} – это сигнал, посылаемый устройством ввода-вывода или процессором (при исключении), который временно переключает CPU на выполнение специального кода – обработчика прерывания.

\subsubsection{Типы прерываний}
В учебнике (главы 2–3) выделяются:
\begin{itemize}
    \item \textbf{Аппаратные прерывания} – генерируются внешними устройствами (диск, клавиатура, таймер) через линии IRQ. Например, прерывание от таймера (IRQ0) происходит 60 раз в секунду.
    \item \textbf{Программные прерывания (ловушки, traps)} – вызываются инструкцией \texttt{int $n$} (например, \texttt{int 0x80} в Linux, \texttt{int SYS386\_VECTOR} в MINIX 3). Используются для системных вызовов.
    \item \textbf{Исключения (exceptions)} – возникают при ошибках в программе (деление на ноль, недопустимая инструкция, нарушение защиты памяти, страничная ошибка). Обрабатываются аналогично прерываниям, но часто приводят к сигналу для процесса.
\end{itemize}

\subsubsection{Для чего нужны прерывания}
\begin{itemize}
    \item Избавляют от необходимости постоянно опрашивать устройства (polling), что экономит ресурсы CPU.
    \item Позволяют CPU переключаться на другие задачи, пока устройство выполняет медленную операцию (например, позиционирование головки диска).
    \item Обеспечивают основу для вытесняющей многозадачности (прерывание от таймера заставляет планировщик переключать процессы).
\end{itemize}

\subsubsection{Как обрабатывается прерывание (на примере Intel 32-bit)}
\begin{enumerate}
    \item Аппаратура сохраняет на стеке текущий программный счётчик (CS:EIP), регистр флагов (EFLAGS) и, в некоторых случаях, SS:ESP.
    \item По номеру прерывания (вектору) из таблицы дескрипторов прерываний (IDT) загружается новый CS:EIP – адрес обработчика.
    \item Обработчик (ассемблерный код) сохраняет все регистры (обычно с помощью \texttt{pushad}) в специальную структуру \texttt{stackframe\_s}, которая находится внутри записи процесса в таблице процессов.
    \item Если прерывание произошло во время выполнения пользовательского процесса, переключается на стек ядра.
    \item Вызывается C-функция \texttt{intr\_handle}, которая определяет, какому драйверу принадлежит прерывание, и отправляет соответствующее уведомление (в MINIX 3).
    \item Планировщик может выбрать другой процесс для выполнения.
    \item По завершении обработки выполняется инструкция \texttt{iretd}, которая восстанавливает регистры из сохранённого контекста и возвращает управление прерванному (или новому) процессу.
\end{enumerate}

\subsection{Защита памяти}
Защита памяти необходима для того, чтобы один процесс не мог читать или изменять память другого процесса или ОС. Она всегда включена для всех процессов, за исключением самых простых встраиваемых систем без MMU (в учебнике такие упоминаются, но MINIX 3 для Intel 32-бит использует защиту).

\subsubsection{Базовый и предельный регистры}
Это не регистры общего назначения, а \textbf{специальные аппаратные регистры}, которые загружаются ОС при переключении контекста.
\begin{itemize}
    \item \textbf{Базовый регистр (base register)} содержит физический адрес начала раздела памяти, выделенного процессу.
    \item \textbf{Предельный регистр (limit register)} содержит размер раздела (в байтах или в единицах, кратных размеру сегмента).
\end{itemize}
Каждый логический адрес, генерируемый процессом, сначала суммируется с содержимым базового регистра, а затем сравнивается с предельным. Если адрес выходит за пределы, генерируется исключение (ошибка сегментации). Такая схема использовалась в CDC 6600 и в упрощённом виде – в Intel 8088 (только база, без предела).

\subsubsection{Ключи и коды защиты (IBM 360)}
В IBM 360 память делилась на блоки по 2 КБ, и каждому блоку присваивался 4-битный \textbf{код защиты}. В регистре состояния программы (PSW) текущего процесса хранился 4-битный \textbf{ключ}. Аппаратура разрешала доступ только если ключ процесса совпадает с кодом защиты блока. Это не криптографические ключи, а простые числовые метки. Операционная система могла менять коды защиты и ключи, а пользовательские процессы – нет.

\subsubsection{Таблица страниц (page table)}
При страничной организации виртуальной памяти каждому процессу соответствует своя таблица страниц (или несколько уровней). Типичная запись таблицы страниц (рис. 4-11) содержит:
\begin{itemize}
    \item \textbf{Номер физического кадра} (page frame number) – основное поле, определяющее, где в физической памяти находится страница.
    \item \textbf{Бит присутствия (present/absent)} – 1, если страница в памяти; 0 – вызывает страничную ошибку (page fault).
    \item \textbf{Биты защиты (protection bits)} – отдельно для чтения, записи и выполнения (или комбинации).
    \item \textbf{Бит модификации (modified/dirty)} – устанавливается аппаратно при записи в страницу; используется при выборе страницы для выгрузки.
    \item \textbf{Бит обращения (referenced)} – устанавливается при любом доступе (чтение или запись); помогает алгоритмам замещения.
    \item \textbf{Бит отключения кэширования} – для страниц, отображающих регистры устройств (memory-mapped I/O).
\end{itemize}
Адрес таблицы страниц хранится в специальном регистре (например, CR3 в Intel). При каждом обращении к памяти MMU извлекает номер виртуальной страницы из адреса, индексирует таблицу и получает номер физического кадра.

\subsection{Преобразование адресов в системах виртуальной памяти}
\subsubsection{Разделение адреса}
Виртуальный адрес делится на \textbf{номер страницы} (старшие биты) и \textbf{смещение} (младшие биты). Размер смещения определяет размер страницы. Например, для 16-битного адреса и страницы 4 КБ (12 бит) номер страницы занимает 4 бита (всего 16 страниц). В 32-битной системе с 4-КБ страницей номер страницы – 20 бит, смещение – 12 бит. Количество бит под номер страницы может быть любым, но обычно степень двойки.

\subsubsection{Физический кадр}
Физический кадр – это блок физической памяти того же размера, что и страница. MMU заменяет виртуальный номер страницы на номер физического кадра, а смещение оставляет без изменений. Полученный адрес отправляется на шину памяти.

\subsubsection{TLB (Translation Lookaside Buffer)}
TLB – это небольшое ассоциативное запоминающее устройство (обычно 8–64 записи), которое хранит наиболее часто используемые отображения виртуальных страниц на физические кадры. Поиск в TLB выполняется параллельно по всем записям. При попадании (hit) физический адрес формируется за один такт, без обращения к таблице страниц в памяти. При промахе (miss) MMU обращается к таблице страниц (обычно через аппаратный обход) и загружает новую запись в TLB, возможно вытесняя старую. В некоторых RISC-архитектурах (SPARC, MIPS, PowerPC) управление TLB выполняется программно: при промахе генерируется исключение, и ОС сама загружает нужную запись.

\subsubsection{Инвертированные таблицы страниц (inverted page table)}
В 64-битных системах виртуальное адресное пространство огромно (например, $2^{64}$ байт). Традиционная таблица страниц с одной записью на каждую виртуальную страницу стала бы астрономического размера. Решение – \textbf{инвертированная таблица}, в которой одна запись соответствует одному физическому кадру, а не виртуальной странице. Запись содержит идентификатор процесса и номер виртуальной страницы, находящейся в этом кадре. Для ускорения поиска используется хеш-таблица, где ключом является (процесс, виртуальная страница). Поскольку в 32-битных системах число виртуальных страниц (около 1 млн) ещё позволяет хранить таблицу в памяти (4–8 МБ), инвертированные таблицы там не применяются – они становятся необходимыми только при переходе на 64 бита, когда прямое отображение уже невозможно.

\subsection{Управление каналами и периферийными устройствами}
\subsubsection{Контроллеры устройств}
Контроллер (адаптер) – это электронная схема (часто отдельная плата), которая управляет физическим устройством и предоставляет программный интерфейс в виде набора регистров. Регистры контроллера могут быть доступны через:
\begin{itemize}
    \item \textbf{Порты ввода-вывода} (специальные инструкции \texttt{IN}, \texttt{OUT}) – используется в x86.
    \item \textbf{Память-отображаемые регистры} (memory-mapped I/O) – регистрам присваиваются адреса в адресном пространстве памяти.
\end{itemize}
ОС записывает команды в эти регистры (например, «прочитать сектор N»), считывает статус и получает данные.

\subsubsection{Модель ввода-вывода (четыре уровня)}
Уровни (рис. 3-8):
\begin{enumerate}
    \item \textbf{Обработчики прерываний} (нижний уровень) – скрывают аппаратные прерывания, преобразуя их в уведомления для драйверов.
    \item \textbf{Драйверы устройств} – код, специфичный для каждого типа устройства. Знает о регистрах контроллера, очередях запросов, обработке ошибок. В MINIX 3 драйверы работают в пользовательском пространстве.
    \item \textbf{Независимый от устройства слой} – выполняет общие функции: буферизацию, кэширование блоков, преобразование имён устройств в номера, проверку прав доступа, выделение/освобождение монопольных устройств.
    \item \textbf{Пользовательский слой} – библиотечные процедуры (printf, scanf), системы спулинга (печать), сетевые демоны.
\end{enumerate}

\subsubsection{DMA (Direct Memory Access)}
DMA – механизм, позволяющий устройству (или специальному контроллеру DMA) передавать данные непосредственно в память без участия процессора. Процессор лишь программирует DMA-контроллер (указывает адрес памяти, количество байт, направление), а затем устройство само пересылает данные по шине. По завершении DMA генерирует прерывание. DMA применяется для дисков, звуковых карт, сетевых адаптеров, когда нужно передавать большие блоки данных. Без DMA процессор вынужден был бы копировать каждый байт из контроллера в память (программируемый ввод-вывод), что загружало бы CPU на 100\% во время пересылки.

\subsection{Firmware — встроенные программы}
В учебнике описаны два типа встроенного ПО:
\begin{itemize}
    \item \textbf{BIOS (Basic Input/Output System)} – находится в ПЗУ на материнской плате. Обеспечивает начальную загрузку (загрузку первого сектора с диска) и простейшие функции ввода-вывода в реальном режиме. MINIX 3 использует BIOS только для первоначальной загрузки, а затем переходит в защищённый режим и работает со своими драйверами.
    \item \textbf{Boot monitor (загрузчик)} – в MINIX 3 это программа \texttt{boot}, которая загружается из первого сектора раздела. Она умеет читать файловую систему, загружать образ ОС, передавать параметры ядру, а также может запускать другие ОС. \textbf{UEFI в книге не упоминается}, следовательно, в ответе следует сказать, что учебник описывает классическую схему BIOS+MBR.
\end{itemize}

\subsection{Middleware — связующее программное обеспечение}
В главе 1 даётся определение: middleware – это ПО, объединяющее распределённые системы и предоставляющее единый интерфейс для приложений (пример – World Wide Web: браузеры и веб-серверы взаимодействуют через HTTP, а промежуточный слой скрывает разнородность).


\subsection{Классификация утилит операционных систем (на примере MINIX 3)}
Утилиты – это программы, не входящие в ядро, но поставляемые вместе с ОС.
\begin{itemize}
    \item \textbf{Файловые утилиты:}
    \begin{itemize}
        \item \texttt{mkfs} – создаёт новую файловую систему на разделе (записывает суперблок, битовые карты, корневой каталог).
        \item \texttt{fsck} – проверяет целостность файловой системы после сбоя (сравнивает счётчики ссылок, ищет потерянные блоки и т.д.).
        \item \texttt{fdisk} / \texttt{part} – управляют таблицей разделов на диске.
    \end{itemize}
    \item \textbf{Системные утилиты:}
    \begin{itemize}
        \item \texttt{ps} – выводит список процессов, их PID, состояние, потребление памяти.
        \item \texttt{shutdown} – корректно завершает работу системы (сбрасывает кэши, посылает сигналы процессам).
        \item \texttt{sync} – принудительно записывает все изменённые блоки из кэша на диск.
        \item \texttt{update} – фоновый процесс, вызывающий \texttt{sync} каждые 30 секунд.
        \item \texttt{chmem} – изменяет размер памяти, выделяемой процессу (поле total в заголовке исполняемого файла).
    \end{itemize}
    \item \textbf{Сетевые утилиты:}
    \begin{itemize}
        \item \texttt{inet} – сетевой сервер (стек протоколов TCP/IP).
        \item Драйверы Ethernet (например, \texttt{rtl8139}) – загружаются отдельно.
        \item \texttt{ftp}, \texttt{telnet} – клиенты для работы по сети (упоминаются как доступные).
    \end{itemize}
    \item \textbf{Утилиты управления устройствами:}
    \begin{itemize}
        \item \texttt{loadkeys} – загружает новую клавиатурную раскладку (ключ-карту).
        \item \texttt{loadfont} – загружает шрифт в видеопамять (для EGA/VGA).
        \item \texttt{stty} – изменяет параметры терминала (скорость, режим, специальные символы).
    \end{itemize}
    \item \textbf{Командный интерпретатор (shell)} – хотя не является частью ОС, выполняет системные вызовы (\texttt{fork}, \texttt{exec}, \texttt{wait}, \texttt{pipe}) и служит основным интерфейсом пользователя.
\end{itemize}

\section*{2. Управление доступом к данным. Файловые системы (основные типы, характеристика)}

\subsection{Управление доступом: модели и механизмы}
В главе 5 подробно разобраны три подхода.

\subsubsection{Матрица доступа (теоретическая модель)}
Строки – домены (процессы), столбцы – объекты (файлы). На пересечении – разрешённые операции (чтение, запись, выполнение). На практике матрица очень разрежена, поэтому её хранят по строкам или столбцам.

\subsubsection{Списки управления доступом (ACL - матрица доступа по столбцам)}
\textbf{Принцип:} у каждого объекта хранится список (домен, права). При доступе система ищет в ACL запись для домена вызывающего процесса.
\begin{itemize}
    \item \textbf{Преимущества:} легко отозвать права (достаточно отредактировать ACL), удобно для объектов с большим числом пользователей.
    \item \textbf{Недостатки:} при каждом доступе может потребоваться поиск по списку (медленнее, чем кейпабилити); проблема с проверкой прав при открытии файла (если права изменились, уже открытые файлы не защищены).
\end{itemize}
В ACL часто используются группы (group), а также дикие карты (wildcard) для всех пользователей. ACL применяются, например, в Windows NT/2000/XP.

\subsubsection{Кейпабилити (capabilities - матрица доступа по строкам)}
\textbf{Принцип:} у каждого процесса есть список (кейпабилити), где для каждого объекта указаны разрешённые операции. При доступе система просто проверяет наличие кейпабилити.
\begin{itemize}
    \item \textbf{Преимущества:} очень быстрая проверка (не нужно искать по ACL), легко реализовать «песочницу» (sandbox) для мобильного кода (дать минимум прав).
    \item \textbf{Недостатки:} сложно отозвать права (нужно найти все кейпабилити, разосланные процессам); требуется защита кейпабилити от подделки.
\end{itemize}
В распределённых системах (Amoeba) кейпабилити защищаются криптографически: сервер хранит случайное число (check field), а клиенту выдаётся хеш от (объект, права, check field). Подделать такое кейпабилити без знания check field невозможно.

\subsubsection{UNIX rwx-биты (очень простая матрица: столбцы - файлы, строки own/group/other)}
Простейший механизм: у каждого файла хранятся 9 бит – по три для владельца, группы и остальных (чтение, запись, выполнение). Дополнительные биты SETUID и SETGID позволяют процессу временно получить права владельца файла.
\begin{itemize}
    \item \textbf{Преимущества:} простота, малый размер, встроен в файловую систему.
    \item \textbf{Недостатки:} грубая градация (только три категории), нет возможности задать разные права для разных групп.
\end{itemize}
В MINIX 3 используется именно эта модель.

\subsection{Типы файловых систем (реализация хранения данных)}
\subsubsection{Непрерывное размещение (contiguous allocation)}
Файл занимает непрерывную последовательность блоков. В каталоге хранятся номер первого блока и длина.
\begin{itemize}
    \item \textbf{Где используется:} CD-ROM, DVD, некоторые старые системы.
    \item \textbf{Достоинства:} высокая скорость последовательного чтения (один seek), простота.
    \item \textbf{Недостатки:} внешняя фрагментация, необходимость заранее знать размер файла.
\end{itemize}

\subsubsection{Связный список (linked list)}
Каждый блок данных содержит указатель на следующий блок. В каталоге – указатель на первый блок.
\begin{itemize}
    \item \textbf{Достоинства:} нет внешней фрагментации.
    \item \textbf{Недостатки:} медленный произвольный доступ (нужно пройти по цепочке), накладные расходы на указатели внутри блоков.
\end{itemize}

\subsubsection{FAT (File Allocation Table)}
Указатели хранятся не в блоках, а в отдельной таблице в памяти. Каталог хранит номер первого блока, а таблица для каждого блока хранит номер следующего (или специальный маркер конца). 
\begin{itemize}
    \item \textbf{Где используется:} MS-DOS, Windows 95/98/ME.
    \item \textbf{Достоинства:} произвольный доступ быстр (вся цепочка в памяти), блоки полностью заняты данными.
    \item \textbf{Недостатки:} таблица FAT занимает много памяти (для 20-ГБ диска с 1-КБ блоками – 80 МБ).
\end{itemize}

\subsubsection{i-узлы (inodes)}
Каждый файл имеет структуру (i-node), содержащую атрибуты и массив указателей на блоки. Обычно 7–12 прямых указателей, затем одно- и двукратная косвенная адресация.
\begin{itemize}
    \item \textbf{Где используется:} UNIX, Linux, MINIX, BSD.
    \item \textbf{Достоинства:} экономия памяти (i-узлы загружаются только для открытых файлов), эффективный произвольный доступ.
    \item \textbf{Недостатки:} для очень больших файлов нужна двукратная (или трёхкратная) косвенная адресация, что увеличивает число обращений к диску.
\end{itemize}
В MINIX 3 i-узел (версии V2/V3) занимает 64 байта, содержит 7 прямых зон и два косвенных указателя (одинарный и двойной). При размере блока 4 КБ максимальный размер файла – более 4 ГБ.

\subsubsection{Журналируемые файловые системы (LFS – Log-Structured File System)}
Все записи (данные, i-узлы, каталоги) буферизуются в памяти, а затем записываются единым большим сегментом в конец журнала (log). Периодически «чистильщик» (cleaner) сканирует журнал и переносит живые данные в начало, освобождая место.
\begin{itemize}
    \item \textbf{Достоинства:} очень высокая производительность при записи маленьких блоков (почти вся запись – последовательная), хорошая масштабируемость.
    \item \textbf{Недостатки:} сложность реализации, накладные расходы на очистку.
\end{itemize}
Применяется в некоторых исследовательских и специализированных ОС (BSD LFS).

\subsubsection{NTFS (Windows NT/2000/XP)}
Файл представляется как набор \textbf{атрибутов}, каждый из которых является потоком байт. Основная структура – MFT (Master File Table), которая содержит до 16 атрибутов. Если атрибут мал, он хранится прямо в MFT (immediate file). Большие атрибуты (например, данные) выносятся в отдельные блоки (nonresident attribute). 
\begin{itemize}
    \item \textbf{Особенности:} Unicode-имена (до 255 символов), пути до 32767 символов, поддержка жёстких и символических ссылок, разрежённые файлы, сжатие, шифрование, журналирование.
    \item \textbf{Достоинства:} высокая надёжность, поддержка больших томов, гибкость.
\end{itemize}

\section*{3. Распределение и использование ресурсов. Алгоритмы планирования}
\subsection{Основные алгоритмы планирования (глава 2)}
\begin{itemize}
    \item \textbf{FCFS (First-Come, First-Served):} процессы выполняются в порядке поступления. Прост, но страдает эффектом «конвоя» (короткие процессы ждут длинных).
    \item \textbf{SJF (Shortest Job First):} выбирается процесс с наименьшим оставшимся временем выполнения. Оптимален по среднему времени ожидания, но требует знания времени выполнения заранее (возможно в пакетных системах).
    \item \textbf{SRTN (Shortest Remaining Time Next):} вытесняющая версия SJF; при появлении нового процесса с меньшим оставшимся временем текущий вытесняется.
    \item \textbf{Round Robin (RR):} каждому процессу даётся квант времени (обычно 20–50 мс). При исчерпании кванта процесс переходит в конец очереди. Хорош для интерактивных систем.
    \item \textbf{Приоритетное планирование:} у каждого процесса есть приоритет; выполняется процесс с наивысшим приоритетом. Может приводить к голоданию низкоприоритетных процессов, поэтому часто применяют динамическое понижение приоритета со временем (aging).
    \item \textbf{Многоуровневые очереди:} процессы делятся на классы (например, интерактивные, пакетные). Внутри каждого класса – свой алгоритм (например, RR для интерактивных, FCFS для пакетных).
    \item \textbf{Гарантированное планирование:} система обещает каждому процессу долю CPU (например, 1/n для n процессов). Выбирается процесс, у которого отношение фактически полученного времени к обещанному минимально.
    \item \textbf{Лотерейное планирование:} процессы получают «лотерейные билеты». Каждое прерывание таймера разыгрывается случайный билет; процесс, владеющий им, получает квант. Доля CPU пропорциональна числу билетов.
    \item \textbf{Fair-share scheduling:} учитывается владелец процессов. Если двум пользователям обещано по 50\% CPU, то независимо от количества их процессов каждому достанется ровно половина.
\end{itemize}
\subsection{Планирование в MINIX 3}
В MINIX 3 используется многоуровневое планирование с 16 очередями (0 – наивысший приоритет, 15 – IDLE). Приоритеты:
\begin{itemize}
    \item Задачи (CLOCK, SYSTEM) – очередь 0.
    \item Драйверы устройств – обычно очереди 1–2.
    \item Серверы (PM, FS) – очереди 3–4.
    \item Пользовательские процессы – по умолчанию очередь 7 (USER\_Q).
\end{itemize}
Процессам даётся квант (для пользовательских – 8 тиков, для системных – 32 или 64). Если процесс использует весь квант и при этом не было других готовых процессов, его приоритет временно понижается (чтобы предотвратить зависание системы). Если процесс блокируется до истечения кванта, при пробуждении он ставится в начало очереди с остатком кванта – это улучшает отклик для ввода-вывода.

\section*{4. Управление памятью. Виртуальная память}
\subsection{Основные методы (без виртуальной памяти)}
\begin{itemize}
    \item \textbf{Монопрограммирование:} вся память занята ОС и одной программой. Просто, но неэффективно.
    \item \textbf{Фиксированные разделы:} память делится на статические разделы при загрузке. Внутренняя фрагментация (неиспользуемое пространство внутри раздела). Использовалось в OS/360 (MFT).
    \item \textbf{Свопинг (swapping):} процессы выгружаются на диск и загружаются обратно. Учёт свободных областей – с помощью битовой карты или списка дыр. Алгоритмы выделения дыр: first fit, next fit, best fit, worst fit, quick fit.
\end{itemize}
\subsection{Виртуальная память: страничная организация}
Виртуальное адресное пространство делится на страницы (обычно 4 КБ). Физическая память – на кадры страниц. Таблица страниц отображает виртуальные страницы в физические кадры. При обращении к отсутствующей странице происходит страничная ошибка (page fault), и ОС подгружает нужную страницу с диска.

\subsection{Алгоритмы замещения страниц}
\begin{itemize}
    \item \textbf{NRU (Not Recently Used):} использует биты R (обращение) и M (модификация). Периодически R сбрасается. При страничной ошибке выбирается случайная страница из наименьшего непустого класса: (0,0), (0,1), (1,0), (1,1).
    \item \textbf{FIFO:} удаляется самая старая страница. Прост, но может вытеснять часто используемые страницы.
    \item \textbf{Second chance:} просматривает R-бит самой старой страницы; если R=0 – удаляет, если R=1 – обнуляет R и переносит страницу в конец очереди.
    \item \textbf{Clock:} циклическая реализация Second chance (часы со стрелкой).
    \item \textbf{LRU (Least Recently Used):} удаляется страница, к которой дольше всего не было обращений. Требует аппаратной поддержки (счётчик или матрица битов). В программной реализации – aging (счётчики сдвигаются, в старший бит добавляется текущий R).
    \item \textbf{Рабочий набор (working set):} множество страниц, использовавшихся за последние k обращений. Если рабочий набор не помещается в память – процесс « thrashing» (свампинг). Алгоритм WSclock использует часы с проверкой принадлежности страницы рабочему набору.
\end{itemize}
\subsection{В MINIX 3}
Виртуальная память не используется. Каждый процесс получает фиксированное количество памяти при fork или exec. Поддерживается разделение текстового сегмента (shared text) для процессов с раздельной I и D памятью.

\section*{5. Организация сетевого взаимодействия}
В главе 1 различаются два подхода:
\begin{itemize}
    \item \textbf{Сетевая ОС:} пользователь знает о существовании нескольких компьютеров, явно выполняет удалённый вход (rlogin), копирует файлы (rcp, ftp). Каждая машина работает под своей локальной ОС, сетевые сервисы добавлены как надстройка.
    \item \textbf{Распределённая ОС:} представляется пользователю как единая система. Распределение ресурсов, миграция процессов, единая файловая система – всё это скрыто от пользователя. Примеры: Amoeba, некоторые исследовательские системы.
\end{itemize}
В MINIX 3 сетевая поддержка реализована отдельным сервером \texttt{inet} (глава 3). Ethernet-драйверы (например, RealTek 8139, Intel Pro/100) работают в пользовательском пространстве. Сервер inet взаимодействует с файловой системой и ядром через сообщения. MINIX 3 может выступать как FTP-сервер или веб-сервер, но распределённой ОС в полном смысле не является.

\section*{6. Виды процессов и управление ими. Средства взаимодействия процессов}
\subsection{Процессы}
Процесс – это программа в стадии выполнения, включающая счётчик команд, регистры, стек, адресное пространство, открытые файлы и т.д. Состояния: running, ready, blocked. Переходы: running $\to$ blocked (при запросе ввода-вывода или семафора), blocked $\to$ ready (по завершении ожидания), running $\to$ ready (по таймеру при вытеснении), ready $\to$ running (планировщик выбирает процесс). Создание: \texttt{fork} (создаёт копию), затем \texttt{exec} (замена образа). Завершение: \texttt{exit}, сигнал, ошибка. Иерархия процессов – дерево, корнем которого является \texttt{init} (PID 1).

\subsection{Потоки (threads)}
Поток – облегчённый процесс, разделяющий с другими потоками того же процесса адресное пространство и ресурсы. У каждого потока свой счётчик команд, стек, регистры. Переключение потоков быстрее переключения процессов (не нужно менять таблицу страниц). Проблемы: синхронизация (race conditions), сигналы (какому потоку доставлять?), управление стеками (переполнение). В MINIX 3 потоки не поддерживаются, но в книге они описаны теоретически.

\subsection{Средства IPC}
\begin{itemize}
    \item \textbf{Семафоры:} целочисленные переменные с атомарными операциями down (P) и up (V). Решают проблему потерянных wakeup. Используются для взаимного исключения (бинарный семафор) и синхронизации (счётный семафор).
    \item \textbf{Мьютексы:} упрощённые бинарные семафоры, удобные для потоков в пользовательском пространстве.
    \item \textbf{Мониторы:} программные модули, в которых компилятор автоматически обеспечивает взаимное исключение. Для ожидания событий – условные переменные с операциями wait и signal. Реализованы в языках Concurrent Pascal, Modula-3, Java (synchronized методы + wait/notify).
    \item \textbf{Обмен сообщениями:} примитивы send и receive. В MINIX 3 используется rendezvous (блокирующий обмен) плюс неблокирующий \texttt{notify} (сохраняется один бит на отправителя). Сообщения фиксированного размера (структура message). Это позволяет строить микроядерные системы, где компоненты ОС общаются через сообщения.
\end{itemize}
Классические задачи: обедающие философы (синхронизация доступа к ресурсам), читатели-писатели (разделение доступа к базе данных).

\section*{7. Модель клиент-сервер. Распределённые ОС}
\subsection{Клиент-сервер в микроядерной ОС}
В микроядерной архитектуре ядро выполняет только самые базовые функции: межпроцессное взаимодействие, низкоуровневое планирование, управление прерываниями. Остальные компоненты (файловая система, драйверы устройств, менеджер процессов, сетевой стек) работают как обычные пользовательские процессы – \textbf{серверы}. Пользовательские программы (\textbf{клиенты}) отправляют серверам сообщения с запросами (например, прочитать файл), сервер выполняет работу и возвращает ответ. Преимущества: модульность (можно заменять серверы независимо), надёжность (сбой драйвера не рушит ядро), простота распределения (сервер может быть на другой машине). MINIX 3 – пример такой системы (PM, FS, RS, inet, драйверы – отдельные процессы).

\subsection{Распределённые ОС}
Распределённая ОС идёт дальше: она скрывает тот факт, что система состоит из множества компьютеров, предоставляя единое виртуальное адресное пространство, единую файловую систему, возможность миграции процессов. Ключевые свойства: прозрачность, масштабируемость, отказоустойчивость. Примеры: Amoeba (разработана Таненбаумом), Mach, некоторые исследовательские системы. В учебнике даётся общее введение, но детальной реализации нет.

\subsection*{Объектно-ориентированный подход в организации ОС}
В учебнике этот подход не рассматривается. Следовательно, в ответе следует указать, что в рамках данного курса (по этому учебнику) вопрос не освещён.

\section*{8. Экспериментальные методы измерения загруженности процессора и использования памяти}
Учебник не содержит систематического изложения экспериментальных методов, но в разных главах имеются следующие фактические сведения:
\begin{itemize}
    \item \textbf{Системный вызов \texttt{times}} (глава 4) – возвращает для процесса: пользовательское время (user time), системное время (system time), а также накопленные времена завершившихся потомков. Это позволяет измерить, сколько CPU времени потребляет процесс.
    \item \textbf{Команда \texttt{ps}} (главы 2 и 4) – отображает текущее состояние процессов: PID, состояние (running, sleeping, zombie), размер памяти, приоритет и т.д. Может использоваться для оценки загрузки памяти в данный момент.
    \item \textbf{Профилирование (profiling)} (глава 2) – если ОС поддерживает системный вызов \texttt{PROF}, драйвер часов может при каждом тике увеличивать счётчик в бине, соответствующем текущему значению программного счётчика. Так строится гистограмма, показывающая, в каких частях программы процесс тратит больше всего времени.
    \item \textbf{Отладочные дампы в MINIX 3} (глава 3) – при нажатии функциональных клавиш (F1–F12) система выводит на экран таблицу процессов (F1), карты памяти (F2), образ загрузки (F3), привилегии процессов (F4), параметры монитора (F5), очереди планировщика (F12) и др. Это позволяет вручную анализировать использование памяти и загрузку CPU.
    \item \textbf{Утилита \texttt{sysenv}} (глава 4) – выводит параметры загрузки, включая блоки памяти, доступные MINIX 3. Показывает общий объём памяти, занятый системой, и свободную память.
\end{itemize}
Детальных методик (например, с использованием счётчиков производительности или имитационного моделирования) в учебнике нет.

\section*{Заключение}
Данный конспект составлен строго по материалам третьего издания учебника Таненбаума и Вудхалла. Все утверждения имеют ссылки на конкретные главы (номера глав указаны в тексте). Вопросы, не рассмотренные в книге (объектно-ориентированный подход, UEFI), выделены отдельно.

\end{document}